\documentclass[aspectratio=169,11pt]{beamer}
\usepackage{../phanes-beamer}
\renewcommand{\ModuleID}{01}
\title{Introduction to agentic engineering}
\subtitle{What language models do, how agents act, and where engineering evidence enters}
\author{PHANES Initiative}
\institute{The University of Texas at Austin}
\date{September 12, 2026}
\begin{document}
\begin{frame}[plain]
\titlepage
\end{frame}
\begin{frame}[t,fragile]{Why introduce agents in nuclear engineering?}
\fontsize{11.5}{14}\selectfont
\begin{itemize}\setlength{\itemsep}{2mm}
\item Engineering work mixes calculations with searches, file preparation, debugging, plotting, and report assembly. Much of the time is spent connecting these steps.
\item An agent can translate a task into proposed edits and tool calls, inspect the returned evidence, and choose a next step.
\item Example: prepare an OpenMC case, diagnose a malformed input, run an approved calculation, and assemble a traceable comparison.
\item Our objective is a reproducible workflow with explicit assumptions, rather than accepting an unverified numerical answer.
\end{itemize}
\end{frame}
\begin{frame}[t,fragile]{AI, language models, and agents are different}
\fontsize{10.5}{12.5}\selectfont\setlength{\tabcolsep}{4pt}\renewcommand{\arraystretch}{1.15}
\begin{tabular}{>{\raggedright\arraybackslash}p{\dimexpr 0.23\linewidth-3.68pt\relax}>{\raggedright\arraybackslash}p{\dimexpr 0.77\linewidth-12.32pt\relax}}
\toprule
\textbf{Term} & \textbf{Working meaning in this course} \\ \midrule
Machine learning & Parameters are fitted from data so a program can make predictions or generate outputs. \\[2pt]
Language model & Predicts a distribution over the next token given the supplied token sequence. \\[2pt]
Token & The smallest unit the model processes: a tokenizer maps text to token identifiers, e.g. 18-bit integers for the Qwen3.8-27B model. \\[2pt]
Assistant & A model interface shaped by instructions, conversation history, and post-training. \\[2pt]
Agentic system & A model plus a harness that can execute tools and continue in response to observations. \\[2pt]
Engineering workflow & A bounded sequence with inputs, methods, acceptance criteria, and retained evidence. \\[2pt]
\bottomrule\end{tabular}\par\vspace{2mm}
\end{frame}
\begin{frame}[t,fragile]{What is a large language model?}
\fontsize{11.5}{14}\selectfont
\begin{itemize}\setlength{\itemsep}{2mm}
\item An LLM is a transformer neural network trained on large text and code corpora; it learns statistical structure, not a project database.
\item Its core operation assigns scores to candidate next tokens. Chat, planning, and tool calls are built around that operation by the surrounding system.
\item It is good at drafting, explaining, rewriting, and producing code or structured output, generalizing from its training distribution.
\item Its limits are known: plausible but wrong details, a knowledge cutoff, no persistent memory of your project, and no access to your systems at inference.
\item For engineering use, treat the model as a probabilistic text generator: every number or citation it emits needs harness or user validation.
\end{itemize}
\end{frame}
\begin{frame}[t,fragile]{The agentic stack has named layers}
\begin{center}
\begin{tikzpicture}[node distance=5mm, every node/.style={draw=UTorange,rounded corners,align=center,font=\scriptsize,minimum height=5mm,minimum width=96mm}, >={Stealth}]
\node (workflow) {\textbf{Engineering workflow}: case files, tools, evidence and acceptance criteria};
\node[below=of workflow] (tools) {\textbf{Tools and skills}: calculations, search, file edits and tests};
\node[below=of tools] (harness) {\textbf{Agent harness}: task loop, permissions and context files};
\node[below=of harness] (api) {\textbf{API endpoint}: OpenAI-compatible interface on the private PHANES node};
\node[below=of api] (weights) {\textbf{Model weights}: qwen3.8-27b-nvfp4 served by vLLM on GPUs};
\begin{scope}[transform canvas={xshift=7mm},thick,->]
\foreach \a/\b in {workflow/tools,tools/harness,harness/api,api/weights}
  \draw (\a.south) -- (\b.north);
\end{scope}
\begin{scope}[transform canvas={xshift=-7mm},thick,->]
\foreach \a/\b in {tools/workflow,harness/tools,api/harness,weights/api}
  \draw (\a.north) -- (\b.south);
\end{scope}
\end{tikzpicture}
\end{center}
\fontsize{10.5}{12.5}\selectfont
\begin{itemize}\setlength{\itemsep}{1.5mm}
\item Requests and responses cross layers in both directions; no single layer owns correctness.
\item Whether the weights and context fit on the GPUs is the budget question; see module 06.
\item OAK harness (\href{https://github.com/jcarlosrodicio/opencode-agent-orchestration-kit}{opencode-agent-orchestration-kit}) implements the same layered pattern for software engineering; module 03 maps its roles in detail.
\end{itemize}
\end{frame}
\begin{frame}[t,fragile]{Open weights versus hosted models}
\fontsize{10.5}{12.5}\selectfont\setlength{\tabcolsep}{4pt}\renewcommand{\arraystretch}{1.15}
\begin{tabular}{>{\raggedright\arraybackslash}p{\dimexpr 0.18\linewidth-4.32pt\relax}>{\raggedright\arraybackslash}p{\dimexpr 0.41\linewidth-9.84pt\relax}>{\raggedright\arraybackslash}p{\dimexpr 0.41\linewidth-9.84pt\relax}}
\toprule
\textbf{Aspect} & \textbf{Open-weight model} & \textbf{Hosted model} \\ \midrule
Access & You download and host the weights; the PHANES node does this. & The provider holds and runs the weights; you send requests. \\[2pt]
Data path & You choose endpoint, logs and retention; module 02 maps the boundary. & Requests traverse provider infrastructure under its stated terms. \\[2pt]
Revision & You pin the served version and update it deliberately. & The model revision may change with provider updates. \\[2pt]
Serving cost & You plan GPU and memory budgets (module 06 works the numbers). & The provider absorbs serving; you pay for access. \\[2pt]
Consequence & The information boundary is explicit, as compliance review needs. & Apply module 02's evidence questions to the provider. \\[2pt]
\bottomrule\end{tabular}\par\vspace{2mm}
\end{frame}
\begin{frame}[t,fragile]{A model generates a conditional sequence}
\[p(x_1,\ldots,x_n)=\prod_{t=1}^{n}p(x_t\mid x_1,\ldots,x_{t-1})\]
\fontsize{11.5}{14}\selectfont
\begin{itemize}\setlength{\itemsep}{2mm}
\item A tokenizer maps text into token identifiers. The model maps a context to scores for possible next tokens.
\item Sampling or decoding selects a token; that token is appended, and generation repeats. A tool call can be one structured output of this process.
\item Training adjusts weights. Inference uses the selected weights; adding a file to a prompt does not retrain them.
\item A plausible continuation can contain a nonexistent API, wrong unit, or unsupported citation. The harness must connect claims to external evidence.
\end{itemize}
\end{frame}
\begin{frame}[t,fragile]{Training stages explain behavior, not authority}
\fontsize{10.5}{12.5}\selectfont\setlength{\tabcolsep}{4pt}\renewcommand{\arraystretch}{1.15}
\begin{tabular}{>{\raggedright\arraybackslash}p{\dimexpr 0.25\linewidth-4.0pt\relax}>{\raggedright\arraybackslash}p{\dimexpr 0.75\linewidth-12.0pt\relax}}
\toprule
\textbf{Stage} & \textbf{Useful distinction} \\ \midrule
Pretraining & Learns statistical structure from a large training corpus. It does not provide a current project database. \\[2pt]
Post-training & Shapes instruction following and other behaviors using additional training signals. \\[2pt]
Inference-time context & Supplies the current question, instructions, evidence, and tool outputs. \\[2pt]
Tool execution & Occurs in the external program environment and may read, edit, compute, or communicate. \\[2pt]
\bottomrule\end{tabular}\par\vspace{2mm}
\end{frame}
\begin{frame}[t,fragile]{Attention and caches explain context costs}
\fontsize{11.5}{14}\selectfont
\begin{itemize}\setlength{\itemsep}{2mm}
\item Attention uses relationships among token representations to build a contextual representation; many models also use sparse or recurrent mechanisms.
\item During autoregressive inference, a runtime can retain intermediate state so each new token does not recompute all preceding work.
\item Long input increases prefill work and consumes runtime state. Long output extends the repeated decode process.
\item The model's advertised context limit, the server's configured limit, and the capacity available under load are three separate quantities.
\end{itemize}
\end{frame}
\begin{frame}[t,fragile]{What is inside one request?}
\fontsize{11.5}{14}\selectfont
\begin{itemize}\setlength{\itemsep}{2mm}
\item A task contains 2,000 instruction tokens, 9,000 history tokens, 13,000 code tokens, and a reserved 4,096-token output allowance.
\item For a 32,768-token configured context, calculate the remaining input margin.
\item A solver returns an 8,000-token log. Decide what to retain verbatim and what to replace with a bounded diagnostic.
\item Explain why dropping the failed-run status while summarizing the log would corrupt the workflow.
\end{itemize}
\end{frame}
\begin{frame}[t,fragile]{Worked solution: the token margin and log triage}
\fontsize{11.5}{14}\selectfont
\begin{itemize}\setlength{\itemsep}{2mm}
\item Remaining margin: $32{,}768 - 2{,}000 - 9{,}000 - 13{,}000 - 4{,}096 = 4{,}672$ tokens before framing overhead.
\item The 8,000-token log exceeds that margin by 3,328 tokens, so it cannot be forwarded verbatim in the next request.
\item Keep verbatim: the failed-run status, the exact command, the case identifier, and the error span. Replace the remainder with a bounded diagnostic plus a path to the full local artifact.
\item Dropping the failure status lets the next agent treat the run as successful and reuse its output; the corrupted quantity is the workflow state, not the log length.
\end{itemize}
\end{frame}
\begin{frame}[t,fragile]{Compare a script, a chat, and an agent}
\fontsize{10.5}{12.5}\selectfont\setlength{\tabcolsep}{4pt}\renewcommand{\arraystretch}{1.15}
\begin{tabular}{>{\raggedright\arraybackslash}p{\dimexpr 0.16\linewidth-3.84pt\relax}>{\raggedright\arraybackslash}p{\dimexpr 0.35\linewidth-8.4pt\relax}>{\raggedright\arraybackslash}p{\dimexpr 0.49\linewidth-11.76pt\relax}}
\toprule
\textbf{Interface} & \textbf{Who chooses the next action?} & \textbf{Engineering consequence} \\ \midrule
Script & Prewritten control flow & Good for defined repetitive calculations; errors need explicit branches. \\[2pt]
Chat & User selects each external action & Useful for explanation; a claimed execution still needs evidence. \\[2pt]
Agent loop & Model proposes; harness applies policy and executes & Can adapt to observations, but requires bounded permissions and validation. \\[2pt]
\bottomrule\end{tabular}\par\vspace{2mm}
\end{frame}
\begin{frame}[t,fragile]{The agent loop has observable states}
\fontsize{11.5}{14}\selectfont
\begin{itemize}\setlength{\itemsep}{2mm}
\item Interpret the task and identify missing inputs; propose a plan with an explicit stopping condition.
\item Request a tool with structured arguments. The harness checks permission and dispatches the operation.
\item Read the tool's status and outputs; revise the plan, retry within a budget, or stop with a diagnostic.
\item Complete only when the requested artifact and acceptance evidence exist. A narrative "done" is not itself a completion condition.
\end{itemize}
\end{frame}
\begin{frame}[t,fragile]{Locate the computation and the data movement}
\fontsize{10.5}{12.5}\selectfont\setlength{\tabcolsep}{4pt}\renewcommand{\arraystretch}{1.15}
\begin{tabular}{>{\raggedright\arraybackslash}p{\dimexpr 0.25\linewidth-4.0pt\relax}>{\raggedright\arraybackslash}p{\dimexpr 0.75\linewidth-12.0pt\relax}}
\toprule
\textbf{Location} & \textbf{What happens there} \\ \midrule
Student Linux host & OpenCode TUI, project files, local tools, Git diffs, tests, and session records. \\[2pt]
Private inference host & vLLM applies model weights to the request context and returns text/tool-call proposals. \\[2pt]
Optional solver host & An explicit local or approved remote command runs the physics code. This path needs its own account and limits. \\[2pt]
Other services & Search, embedding, MCP, plugins, or file synchronization can create additional destinations. \\[2pt]
\bottomrule\end{tabular}\par\vspace{2mm}
\end{frame}
\begin{frame}[t,fragile]{A structured call separates intent from execution}
\fontsize{11.5}{14}\selectfont
\begin{itemize}\setlength{\itemsep}{2mm}
\item Model proposal: invoke a calculation tool with named SI inputs; the harness validates the request shape and permission.
\item Execution result: process exit status, parsed JSON, raw-output path, and the case identifier.
\item Application validation: check units, finiteness, completeness, and correspondence to the approved input.
\item Only the validated result can enter the final report. A string that looks like a shell command is not proof it was executed.
\end{itemize}
\end{frame}
\begin{frame}[t,fragile]{Inspect a concrete tool contract}
\begin{lstlisting}
INPUT:  power_W=100000, mass_flow_kg_s=2,
        cp_J_kg_K=1000, inlet_K=300
OUTPUT: {"status":"ok", "delta_T_K":50.0,
         "outlet_K":350.0, "units":{"outlet_K":"K"}}
FAILURE: nonzero exit + explicit error; no new success file
\end{lstlisting}
\fontsize{10.5}{12.5}\selectfont
\begin{itemize}\setlength{\itemsep}{2mm}
\item The reference equation is $T_{\mathrm{out}} = T_{\mathrm{in}} + \frac{Q}{\dot{m} c_p}$ (in script: \texttt{inlet\_K + power\_W / (mass\_flow\_kg\_s * cp\_J\_kg\_K)}).
\item The model can propose the run and explain it; Python performs this calculation.
\item Check whether changing mass flow causes a fresh execution with a distinct result file.
\end{itemize}
\end{frame}
\begin{frame}[t,fragile]{Context files have different jobs}
\fontsize{10.5}{12.5}\selectfont\setlength{\tabcolsep}{4pt}\renewcommand{\arraystretch}{1.15}
\begin{tabular}{>{\raggedright\arraybackslash}p{\dimexpr 0.25\linewidth-4.0pt\relax}>{\raggedright\arraybackslash}p{\dimexpr 0.75\linewidth-12.0pt\relax}}
\toprule
\textbf{Artifact} & \textbf{What belongs in it} \\ \midrule
Task request & Current objective, scope, desired artifacts, and acceptance conditions. \\[2pt]
AGENTS.md & Persistent project conventions: commands, units, file ownership, and required checks. \\[2pt]
Skill instructions & A reusable procedure with triggers, inputs, outputs, references, and stop conditions. \\[2pt]
Run manifest & Actual inputs, versions, hashes, process status, and output locations. \\[2pt]
Documentation & How another engineer configures, runs, checks, and interprets the workflow. \\[2pt]
\bottomrule\end{tabular}\par\vspace{2mm}
\end{frame}
\begin{frame}[t,fragile]{Reasoning, retrieval, and tools solve different problems}
\fontsize{11.5}{14}\selectfont
\begin{itemize}\setlength{\itemsep}{2mm}
\item Reasoning helps organize a difficult task, but an extended explanation can still rest on a false premise.
\item Retrieval supplies selected external text. Verify its source, revision, and applicability before treating it as evidence.
\item Tools provide observations or execute operations. Their contracts and implementation can also contain defects.
\item Combine all three with independent checks: a correct source, a validated input, an actual run, and a reviewed interpretation.
\end{itemize}
\end{frame}
\begin{frame}[t,fragile]{Configure the private endpoint explicitly}
\begin{lstlisting}
PHANES_BASE_URL='https://inference.example.edu:41883/v1'
PHANES_API_KEY='replace-with-your-issued-key'
PHANES_MODEL='qwen3.8-27b-nvfp4'
PHANES_CONTEXT=262144
PHANES_OUTPUT=8192
\end{lstlisting}
\fontsize{10.5}{12.5}\selectfont
\begin{itemize}\setlength{\itemsep}{2mm}
\item The trusted shared/.env file is explicitly sourced by start-opencode.sh.
\item The wrapper generates a provider entry whose model key matches PHANES\_MODEL.
\item The credentials are placeholders; use the issued endpoint and keep the real key out of submissions.
\end{itemize}
\vfill{\fontsize{6.5}{8}\selectfont\color{UTgray} Sources: occonfig, user. See references and instructor notes.\par}
\end{frame}
\begin{frame}[t,fragile]{Use the compatible API adapter}
\begin{lstlisting}
"provider": {"phanes": {
  "npm": "@ai-sdk/openai-compatible",
  "options": {"baseURL": "{env:PHANES_BASE_URL}",
              "apiKey": "{env:PHANES_API_KEY}"},
  "models": {"qwen3.8-27b-nvfp4": {}}
}}
\end{lstlisting}
\fontsize{10.5}{12.5}\selectfont
\begin{itemize}\setlength{\itemsep}{2mm}
\item This is a configuration excerpt. Use shared/start-opencode.sh for the complete generated file.
\item The URL includes /v1; baseURL is case-sensitive, and the model identifier must match the served alias.
\end{itemize}
\vfill{\fontsize{6.5}{8}\selectfont\color{UTgray} Sources: ocprovider. See references and instructor notes.\par}
\end{frame}
\begin{frame}[t,fragile]{Run the first controlled TUI session}
\begin{lstlisting}
cd shared
cp .env.example .env       # edit the issued settings
chmod 600 .env
./start-opencode.sh --check
./start-opencode.sh
\end{lstlisting}
\fontsize{10.5}{12.5}\selectfont
\begin{itemize}\setlength{\itemsep}{2mm}
\item In /models, confirm the phanes provider and expected alias.
\item Ask for a plan and a local calculator run; inspect the proposed command before approving it.
\item Read the saved JSON, compare with the analytical equation, and inspect the final explanation.
\end{itemize}
\end{frame}
\begin{frame}[t,fragile]{Diagnose failures by layer}
\fontsize{10.5}{12.5}\selectfont\setlength{\tabcolsep}{4pt}\renewcommand{\arraystretch}{1.15}
\begin{tabular}{>{\raggedright\arraybackslash}p{\dimexpr 0.25\linewidth-4.0pt\relax}>{\raggedright\arraybackslash}p{\dimexpr 0.75\linewidth-12.0pt\relax}}
\toprule
\textbf{Observation} & \textbf{First evidence to inspect} \\ \midrule
TLS error & Hostname, certificate chain, clock and institutional CA configuration. \\[2pt]
HTTP 401 & Issued key, authorization header and server-side account policy. \\[2pt]
Unknown model & /v1/models versus the generated provider's model-map key. \\[2pt]
No tool call & Harness trace, model capabilities, server parser and chat template. \\[2pt]
Wrong numerical answer & Actual command arguments, tool version, units and validated output. \\[2pt]
\bottomrule\end{tabular}\par\vspace{2mm}
\end{frame}
\begin{frame}[t,fragile]{What the user contributes during development}
\fontsize{11.5}{14}\selectfont
\begin{itemize}\setlength{\itemsep}{2mm}
\item At the start: a bounded question, source data, units, constraints, and a definition of done.
\item At a decision: resolve a consequential ambiguity or authorize a specific operation with visible effects.
\item During execution: inspect progress artifacts and failures; correct the objective if the work drifts.
\item At completion: review the diff, tests, provenance, limitations, and final result before accepting it.
\end{itemize}
\end{frame}
\begin{frame}[t,fragile]{Improve this engineering request}
\fontsize{11.5}{14}\selectfont
\begin{itemize}\setlength{\itemsep}{2mm}
\item Weak request: "Make a reactor model and tell me if it works."
\item Rewrite it with a public teaching geometry, specified outputs, a compute budget, and explicit validation stages.
\item Require the agent to list unresolved physics assumptions before implementing them.
\item Require both a successful artifact and a deliberate failure demonstration.
\end{itemize}
\end{frame}
\begin{frame}[t,fragile]{One worked approach: improving the engineering request}
\fontsize{11.5}{14}\selectfont
\begin{itemize}\setlength{\itemsep}{2mm}
\item Rewritten request: implement the provided public pin-cell specification; export the OpenMC XML and a geometry plot; validate both against the specification; run only the approved smoke budget.
\item The agent lists unresolved physics assumptions before implementing them and retains the diagnostics.
\item It demonstrates one controlled failure with an invalid input and confirms that no success artifact is written.
\item It requests approval before any convergence study; no physical conclusion is claimed from a smoke run.
\end{itemize}
\end{frame}
\begin{frame}[t,fragile]{Lab deliverable: explain the whole evidence chain}
\fontsize{11.5}{14}\selectfont
\begin{itemize}\setlength{\itemsep}{2mm}
\item Capture the raw API response, local tool invocation, saved calculation output, and final TUI explanation.
\item Annotate which component supplied each capability: model, harness, tool, validator, or user.
\item Include one invalid-input case and one endpoint-configuration diagnosis.
\item Submit a short explanation of what crossed the inference boundary and what evidence supports the reported value.
\end{itemize}
\end{frame}
\begin{frame}[t,fragile]{Sample submission: the evidence-chain deliverable}
\fontsize{11.5}{14}\selectfont
\begin{itemize}\setlength{\itemsep}{2mm}
\item Four captured artifacts: the raw API response, the exact local tool command, the saved JSON output, and the final TUI explanation.
\item Annotation table: the model supplied the plan and wording; the harness enforced permissions and dispatch; the tool performed the calculation; the validator checked units and finiteness; the user authorized the endpoint and the run.
\item Invalid-input case: a negative mass flow produced a nonzero exit and an explicit error, with no new success file.
\item Endpoint diagnosis: an HTTP 401 was traced to the issued key and authorization header, not to the model settings.
\item Boundary statement: the request context crossed to the private endpoint; the reported 350.0 K is supported by the saved JSON and the analytical equation, not by prose.
\end{itemize}
\end{frame}
\begin{frame}[t,fragile]{Before the next module, answer four questions}
\fontsize{11.5}{14}\selectfont
\begin{itemize}\setlength{\itemsep}{2mm}
\item What makes this workflow agentic rather than a fixed script?
\item Where did the numerical work actually run, and how can you prove it?
\item Which files or tool outputs became model context?
\item Which changes to data, recipients, tools, or endpoints require revisiting the information boundary?
\end{itemize}
\end{frame}
\begin{frame}[t,fragile]{A complete answer: the four pre-module questions}
\fontsize{11.5}{14}\selectfont
\begin{itemize}\setlength{\itemsep}{2mm}
\item Agentic: the model proposes the next action from observed tool results while the harness enforces policy; a fixed script has prewritten control flow.
\item The numerical work ran in the local Python process; the proof is the exit status, the saved JSON with its file path, and the run manifest. The endpoint returned text only.
\item Model context included the task text, project context files, and the tool's structured output inserted by the harness.
\item Revisit the information boundary when the data class, the recipient service, a network-capable tool, or the endpoint or provider changes.
\end{itemize}
\end{frame}
\begin{frame}[t]{References 1}
\fontsize{9}{12}\selectfont\begin{itemize}\setlength{\itemsep}{3mm}
\item \textbf{curriculum}: User-supplied PHANES curriculum: mod/01 through mod/08/README.md.
\item \textbf{colors}: \href{https://umac.utexas.edu/brand-center/colors/}{umac.utexas.edu/brand-center/colors}
\item \textbf{occonfig}: \href{https://opencode.ai/docs/config/}{opencode.ai/docs/config}
\item \textbf{user}: User-supplied vllm(1).zip: vllm/vllm-qwen38-nvfp4.sh, README.md, chat\_template.jinja; reviewed 2026-09-05.
\item \textbf{ocprovider}: \href{https://opencode.ai/docs/providers/}{opencode.ai/docs/providers}
\item \textbf{oak}: \href{https://github.com/jcarlosrodicio/opencode-agent-orchestration-kit}{github.com/jcarlosrodicio/opencode-agent-orchestration-kit}
\end{itemize}\end{frame}
\end{document}
